\begin{abstract}
Which channels should be kept when a pretrained network is pruned to recognise classes it was never trained on, from a handful of labelled examples? We compare eight channel scores on a ResNet-18 and five of them on a MobileNetV2: random choice, filter $\ell_1$ norm, BatchNorm scale, a first-order Taylor term, DACIS (a rank sum of three statistics) and each of those statistics alone. Test classes are held out from NWPU-RESISC45 aerial scenes and 13-band Sentinel-2 EuroSAT-MS patches, with leaf photographs as a ground-level control; five to ten random class partitions, one set of test episodes for all methods. DACIS is never reliably ahead of random choice. The Taylor term is ahead of it by \key{\num{d:rs:novel_5:taylor-random:PMP:abs}} points on RESISC45 for the ResNet-18 at $30\%$ of the parameters, and by \key{\num{d:rs:novel_5:taylor-random:PMP10:abs}} and \key{\num{d:rs:novel_5:mp-taylor-random:25:abs}} points at the severest budgets we tried (ResNet-18 at $10\%$, MobileNetV2 at $25\%$), where BatchNorm scale on the MobileNetV2 falls \key{\num{d:rs:novel_5:mp-bn_gamma-random:25:abs}} points below random choice. The meta-learning stage of a prune, meta-train, prune schedule did not raise 5-shot accuracy, and a naive 13-band input did not beat RGB. Pruning the ResNet-18 to $30\%$ removed \key{\num{cost:macs-red:30}\%} of the multiply-accumulates but only \key{\num{cost:cpu-red:30}\%} of CPU latency, and an unpruned MobileNetV3-Small, at \num{cost:mobilenet_v3_small:macs} G multiply-accumulates, beats every pruned network we built on the three novel-class sets. Notebooks rerun the experiments and rebuild every table.
\end{abstract}
